Triangle $BKL$ is isosceles with vertex angle $52^\circ$, so $\angle BKL=64^\circ$. Since $KA$ and $KB$ are opposite rays, $\angle AKL=116^\circ$.
Triangle $AKC$ is isosceles with $AK=KC$, so $\angle KAC=\angle ACK=40^\circ$, giving $\angle AKC=100^\circ$. Both rays $KL,KC$ lie in the same half-plane above $AB$; subtracting gives
\[
\boxed{\angle LKC=116^\circ-100^\circ=16^\circ}.
\]
\begin{asy}
texpreamble("\usepackage{amsmath,amssymb}");
size(240);
defaultpen(fontsize(9pt));
pair A=(0,0),C=dir(40),B=(Sin(88)/Sin(52),0),K=(1/(2*Cos(40)),0),L=B+abs(B-K)*(C-B)/abs(C-B);draw(A--B--C--cycle);draw(K--C,blue);draw(K--L,red);label("$A$",A,(-1,-1));label("$B$",B,(1,-1));label("$C$",C,(0,1));label("$K$",K,(0,-1));label("$L$",L,(1,0));
\end{asy}
